\section{Methodology}
\label{sec:methodology}

This section describes and analyses in detail the procedure and process for evaluating the AI tools described (Windsurf Cascade and Github Copilot). The aim of the study was to evaluate current AI tools, such as Windsurf Cascade and Github Copilot in the context of Infrastructure as a Service, more specifically Terraform, in order to be able to make statements about how successfully such AI-supported programming aids can be used in the real world and how reliable and robust the generated code is. In addition, the aim was to examine how current AI programming tools in the IaC area deal with tasks of varying degrees of difficulty and, if possible, to determine the areas in which these tools can currently be used most effectively. It also sought to identify which aspects of the tools need to be further improved and researched so that they can be classified as truly reliable and trustworthy in code generation in the IaC area.

The evaluation was carried out using the two AI tools already mentioned: Windsurf Cascade was evaluated using the Windsurf code editor and Github Copilot was evaluated using Visual Studio Code. In order to evaluate both systems under as similar conditions as possible, the Claude 4 Sonnet model from Antropic \cite{claude_sonnet_2025} was selected as the standard AI model. It is important to note that this evaluation is not about evaluating the model, but only about evaluating the AI systems. Both editors and systems were configured so that no previous chats or memory data were included in the execution and results of the system.

A total of 11 prompts from three categories (easy, medium, complex) will be created with a distribution of 3, 4, 4 and briefly evaluated by a pilot test on both AI systems to check the wording of the prompts and ensure that there is still enough leeway for the model to interpret the prompt independently. The tasks range from creating simple S3 buckets to very complex ones such as a 3-tier architecture with load balancer, web tier and application tier. An example of a simple prompt looks like this: ''Create a single EC2 t3.micro instance with SSH security group allowing access from anywhere in us-east-1 region''. A prompt from the complex category looks like this: ''Create EKS cluster with managed node group using t3.medium instances, including all required IAM roles, policies, and security groups''. After the pilot test, a fixed list of the 11 prompts mentioned will be compiled, which can be found in the attached Github repository. 

The evaluation process for AI systems is as follows: An empty folder is created in a test environment and opened using the respective editor, the AI programming assistant is opened and configured to the Claude 4 Sonnet model, and the editing mode is set to ''edit'' in Github Copilot or ''chat'' in Windsurf Cascade. The system is checked to ensure that it has no saved memories or previous chats; if any are found, they are deleted. The respective prompt is then inserted from the file containing the defined prompts into the input field of the AI programming assistant. Next, a ''main.tf'' file is created and opened. This file ensures that all of the system's code is focused on a single file so that the system's handling of complex scenarios without subdivisions can be evaluated simultaneously. After the file is created, the current time is noted, a timer is started, and the prompt is sent to the system. The system waits until the prompt has been completely processed and indicates that it is ready. At this point, the timer is stopped and the current time is noted again. To evaluate the code written by the system, the code is permanently stored in the file using the ''Apply'' function in Windsurf Cascade or the ''Keep'' function in Github Copilot. To save the code stored in the file, it is copied to a file in the repository found in the appendix. To evaluate the code and check it for functionality and possible errors, the commands ''terraform init'', ''terraform validate'' and ''terraform plan'' are executed in the current editor. The respective results of these commands are then also stored in separate files in the repository found in the appendix. To complete the respective run, the code is checked and evaluated for functionality, quality and completeness. The evaluation scheme comprises 0-2 points per parameter, whereby decimal numbers are also possible. The results of these evaluations are then added up and stored in an Excel file along with all other data, such as the time, start and end times, AI assistant, prompt ID, run number, comments and a description of whether the execution of the ''keep'' or ''apply'' function was successful or whether the code had to be inserted manually. This can also be found in the repository listed in the appendix. To ensure the quality and transparency of the individual steps, each of these steps is documented with a screenshot of the current status and also stored in the repository. 

To make the results more meaningful and ensure that the quality and performance of the systems are thoroughly analyzed, each prompt is executed three times on each system according to the procedure described above. With 11 prompts, two systems and 3 executions per prompt and system, this results in a total of 66 executions (33 Windsurf Cascade and 33 GitHub Copilot via VSCode). The respective runs are performed depending on the prompt, i.e. first 3 executions of prompt ‘E1’ are performed via Github Copilot and then 3 executions of the same prompt via Windsurf Cascade. Then the system switches to the next prompt, in this case prompt ‘E2’, and the same procedure is carried out systematically until all prompts have been processed one after the other and all necessary data has been saved in the Excel table mentioned above. The executions are given a time limit of max. 5 minutes per execution and prompt. In addition, all exact configurations and versions of Terraform, Copilot, Windsurf and all other tools used can also be found in a file in the aforementioned repository.

In order to be able to make actual statements after the implementation and evaluation of the systems about how well the respective systems performed with different levels of complexity of prompts, in general and also in comparison, a detailed analysis of the results is required. This is in order to determine what assessments can be made about the systems and their use cases. For this purpose, the data entered in the Excel file mentioned above is analyzed and evaluated in a structured manner. In order to analyze and evaluate the amount of data as simply and efficiently as possible, a Python script is written that takes the Excel file as input and reads, saves and then evaluates the data from the created table. The following data is to be analyzed: Speed of the systems evaluated individually and against each other in relation to all tasks as well as task groups such as only complex or only simple tasks. In addition, the success rate of successful validation and planning is to be evaluated, i.e. the results of executing the commands ‘terraform validate’ and ‘terraform plan’. And, of course, data such as the general point score of the systems and the point score for individual categories with the resulting standard deviation. In addition, individual categories such as quality or functionality as a whole, but also categories related to the tasks, should be evaluated.

In order to present the data obtained in a clearer manner, Python and the matplotlib software package will be used to generate graphics that clearly show the most important results of the evaluation, thus providing a simple overview of the evaluation results and enabling users to make their own assessment of the relationship between the systems. Graphics such as heat maps, distributions and simple bar charts are used to present the results in a simple and meaningful way.